\scanpage{42}
Appelons F-V-\uline{\unsure{cristal}} \uline{filtré}\inserted{ \uline{sur} S} (par abus de langage) un F-V-cristal en Mod. loc. lib. sur S muni de\deleted{une filtr} la donnée d'une filtration de $M_S$ par une sous-module localement facteur direct. On trouve donc en fait un foncteur
\[
\diagram{\begin{tikzcd}
\mathrm{BT}(S) \arrow[r] & \text{F-V-cris fil }(S),
\end{tikzcd}}
\]
compatible aux images inverses.

\inserted{\uline{Remarque}}
La filtration (*) doit être considérée comme étant l'équivalent d'une filtration ``de Hodge'' sur une \unsure{cohomoologie|cohomologie} de De Rham \inserted{relative} en dimension 1. \deleted{Pour expliciter} Ce point\deleted{, partons par exemple \unsure{d'un schéma abélien} A sur S, et prenons $G=\unsure{A}(\infty)=\unsure{\lim|\varinjlim}A(n)$. On peut alors}\inserted{ sera explicité plus bas.}

Soit maintenant S' un épaississement à p.div. de S. On se propose de trouver tous les prolongements possibles (à isom. près) \unsure{da|du} groupe de BT G donné sur S en un groupe de BT G' sur S'. On \deleted{trouve} prouvera que si à tout tel G', on associe la filtration correspondante de $D(G')_{S'}=\underline D^*(G)_{S'}=\underline D^*(G_0)_{S'}$ (cf.\qquad pour le fait que S' peut être envisagé comme un objet du site cristallin absolu de $S_0$), qui prolonge la filtration qu'on a sur $\underline D^*(G)_S=\underline D^*(G_0)_S=\underline D^*(G')|S$, on trouve une bijection entre l'ensemble des classes d'isomorphie cherchés et l'ensemble des filtrations qui prolongent. De façon plus précise, on a:\deleted{une équivalence de catégories}

\margin{\figure{\begin{tikzpicture}\draw (0,0)--(0,1.5);\end{tikzpicture}}}
\uline{Théorème de déformation pour les groupes de} BT : Soient S un schéma où p est loc. nilp., S' un épaississement à p.div. de S, Considérons le foncteur canonique
\[
\diagram{\begin{tikzcd}[column sep=large]
\mathrm{BT}(S') \arrow[r] & \begin{array}{l}
\text{catégorie des }\deleted{\text{triples}}\inserted{\text{paires}}\text{ d'un groupe de BT G sur S, et}\\
\text{d'un}\deleted{\text{e}}\text{ sous-module localement facteur direct de}\\
\underline D^*(G)_{S'}\text{ qui prolonge le sous-module }\operatorname{Fil}^1\underline D^*(G)_S=\underline\omega_G\\
\text{de }\underline D^*(G)_S.
\end{array}
\end{tikzcd}}
\]
\margin{\figure{\begin{tikzpicture}\draw (0,0)--(0,.6);\end{tikzpicture}}}
Ce foncteur est une \uline{équivalence de catégories}, \inserted{\uline{si} S' \uline{est à p. div. nil} ou si on se borne à des G qui sont infinitésimaux, ou à \unsure{duaux} infinitésimaux.}
\figure{\begin{tikzpicture}\draw (0,.1)..controls (3,.45) and (5,-.15)..(7,.1)--(13,.1);\end{tikzpicture}}

\inserted{Remarques}
Si le \deleted{\unsure{questio|question}} problème 1 avait une réponse affirmative, on en conclurait une description de la catégorie des groupes de BT sur S en termes de cristaux de Dieudonné filtrés sur S, en appliquant le théorème de
\margin{41}
